\subsection{Modular Flavour Symmetries and Modulus Stabilisation --- arXiv:2201.02020}

\textbf{Authors:} P. P. Novichkov, J. T. Penedo, S. T. Petcov

\paragraph{主な主張}
モジュラー対称性と CP を保つ単純な超重力由来ポテンシャルに、$\tau=\omega=e^{2\pi i/3}$ の近傍で CP を自発的に破る極小が微調整なしに存在することを示した \cite{Novichkov:2022wvg}。

\paragraph{新規性}
固定点からの小さな変位と CP 破れを同時に動的に生じさせ、フェルミオン質量階層に必要な小パラメータの起源を与えた。

\paragraph{位置付け}
modular flavor 模型で入力とされがちなモジュラスの値を、残留対称性近傍の真空構造から説明する基礎研究である。

\paragraph{コメント（読書メモ）}
$\tau=\omega$ 近傍での固定化を扱う。読書時点のメモでは、dS 真空は見つけられていないと整理している。
